\section{Description Logics with Concrete Domains}
\label{sec:alcd}

Here, we present the logic $\alcd$~\cite{BaaderH91}.
This is the logic $\alc$ extended with:
  \emph{abstract} and \emph{concrete features}, and, more importantly, a
  \emph{concrete domain}~$\cdom$.

\begin{definition}[Concrete Domain]\label{def:conc:domain}
    A \emph{concrete domain} is a pair $\cdom = \tup{\Delta_\cdom, \Phi_\cdom}$ where:
        $\Delta_\cdom$ is a non-empty set, and
        $\Phi_\cdom$ is a set of $n$-ary predicate names s.t.\
        each $P \in \Phi_\cdom$ has an extension $P_\cdom \subseteq \Delta_\cdom^n$.
    A concrete domain is \emph{admissible} if:
        (i) $\Phi_\cdom$ it is closed under negation:
            for each $P\in\Phi_\cdom$ there is $\bar{P}\in\Phi_\cdom$ such that $\bar{P}^\cdom=\Delta_\cdom\setminus P_\cdom$,
        (ii) exists $\toppred \in \Phi_\cdom$ s.t.\ $\toppred_\cdom=\Delta_\cdom$, and
        (iii) exists an algorithm for deciding satisfiability of finite conjunctions of the form
            $\bigwedge\setof{P_i(\mathbf{x}_i)}{ i \in [1 \dots k]}$, where $P_i \in \Phi_{\cdom}$ and $\mathbf{x}_i$ is a sequence of variables whose length is the arity of $P_i$.
\end{definition}

We introduce the syntax and semantics of $\alcd$ relative to an arbitrary but fixed concrete domain $\cdom$ below.
To this end, let
  $N_C$ be a countable set of \emph{concept names},
  $N_R$ be a countable set of \emph{role names},
  $N_{\mathit{aF}} \subseteq N_R$ be a countable set of \emph{abstract features}, and
  $N_{\mathit{cF}}$ be a countable set of \emph{concrete features}. 
The elements in $N_{\mathit{aF}}^{*}$ are called \emph{abstract paths}.
% Since our goal is to use concrete domains to reason about $\xpath$, we consider a simple concrete domain $\cdom = \set{\D, \{= , \neq, \toppred\} }$.
% It turns out that embedding $\xpath$ into $\alcd$ over this particular concrete domain will give us interesting complexity bounds. More interestingly, we will refine these complexity bounds and will see $\xpath$ as a fragment of $\alcd$ with good computational properties. 

  
% \begin{definition} [Signature] \hfill
%     \begin{enumerate}
%       \item $N_C=\set{C_1, C_2 , \dots}$ is a contably infinite set, called \emph{concept names};
%       \item $N_R= \set{r_1, r_2 , \dots}$  is a contably infinite set, called  \emph{role names};
%       \item $N_{aF}=\set{f_1, f_2 , \dots}$ is a  infinite subset of $N_R$,  called \emph{abstract features};
%       \item $N_{cF}=\set{g_1, g_2 , \dots}$ is a infinite set, called \emph{concrete features}.
%       \item $O_a= \set{a_1, a_2, \dots}$ is a contably infinite set called \emph{abstract object};
%       \item $O_c= \set{x_1, x_2, \dots}$ is a contably infinite set called \emph{abstract object};
%     \end{enumerate}
% \end{definition}

\begin{definition}[Language] \label{def:alcd-syntax}
  % Let $\cdom$ be a concrete domain. % as in \Cref{def:conc:domain}.
    The $\alcd$ \emph{concepts} are defined by the grammar:
    \[
        C, D \Coloneqq
        C_i \mid
        \lnot C \mid
        C \sqcap D \mid
        %C \sqcup C \mid 
        \exists R.C \mid
        \exists u_1g_1 \dots u_ng_n.P,
    \]
    where $C_i\in N_C$, $R\in N_R$, $P \in \Phi_\cdom$ is a $n$-ary predicate name, each $u_i$ is an abstract path, and $g_i \in N_{\mathit{cF}}$.
\end{definition}

The semantics of $\alcd$ concepts is given in terms of \emph{interpretations} and concrete domains. Later on we will relate data models, and interpretations and concrete domains as alternative presentations of the same kind of structure.

\begin{definition}[Interpretation]\label{def:interpretation}
    An \emph{interpretation} is a pair
    ${\interp = 
    \tup{\Delta_{\interp}, {\cdot^{\interp}}}} $ 
    where:
    \begin{enumerate}
      \item 
        $\Delta^{\interp}$ is a non-empty set of elements (also called domain).
      \item 
        $\cdot^{\interp}$ is a function that maps:
        \begin{enumerate}
          \item 
            every $C_i \in N_C$ to a subset $\sem{C_i} \subseteq \Delta_{\interp}$;
          \item  
            every $r_i \in N_R$ to a binary relation $\sem{r_i}$ on $\Delta_{\interp}$;
          \item 
            every  $f_i \in N_{aF}$, to a partial function $\sem{r_i}$ on $\Delta_{\interp}$;
          \item 
            every $g_i \in N_{cF}$, to a partial function $g_i^\interp : \Delta_\interp \nrightarrow \Delta_D$.
          % \item 
          %   every $a_i \in O_a$, $\sem{a_i} \in \Delta_{\interp}$;
          % \item 
          %   every $x_i \in O_c$, $\sem{x_i} \in \Delta_{D}$;  
        \end{enumerate}
    \end{enumerate}
    For every abstract path $u = f_1 \dots f_k$, the interpretation $u^{\interp}$ of $u$ is defined as the (relational) composition $f_1^\interp; \dots; f_k^\interp$.
    % Moreover, if $g$ is a concrete feature, the interpretation $(ug)^{\interp}$ of $ug$ is defined as the (relational) composition $u^\interp; g^\interp$.
  \end{definition}
  
  \begin{definition} [Semantics] \label{def:alcd:semantics} 
    Let ${\interp = \tup{\Delta_{\interp}, {\cdot^{\interp}}}}$ be an interpretation, and $\cdom = \tup{\Delta_\cdom, \Phi_\cdom}$ be a concrete domain.
    The semantics of $\alcd$ concepts is: 
    \begin{align*}
      (\lnot C)^\interp                  ={}& \Delta^\interp \setminus C^\interp \\
      (C \sqcap D )^\interp              ={}&  {C^\interp} \cap {D^\interp} \\
      (\exists R.C )^\interp             ={}& \setof{x}{ \text{there exists } y \text{ such that } (x, y) \in R^\interp \text{ and } y \in C^\interp } \\
      (\exists u_1g_1 \dots u_ng_n.P)^\interp  ={}&  \{x \mid \text{there exists } {y_i \in \dom(g_i^\interp)} \text{ such that } \\
      & \qquad (x,y_i) \in u_i^\interp \text{ and } (g_1^\interp(y_1) \dots g_n^\interp(y_n)) \in P_\cdom \}.
  \end{align*}
  \end{definition}


  The results listed below are relevant to provide some complexity bounds of reasoning tasks for $\xpath$ in a DL setting.  See~\cite{Lutz1998-LUTDLW-2} for further details.

  \begin{proposition}[\cite{Lutz02}] \label{prop:lutz}
    If $\cdom$-satisfiability is in \pspace, then we have the following complexity bounds for $\alcd$ reasoning problems: 
    \begin{enumerate}
      \item concept satisfiability is \pspace-complete;
      \item ABox consistency is \pspace-complete;
      \item concept satisfiability w.r.t.\ acyclic TBoxes is in \nexptime.
    \end{enumerate}
  \end{proposition}
  